\chapter{配置、复位与统计}\label{ch:cfg}

\code{afdx_cfg} 负责配置总线的地址译码、配置寄存器和 27 个统计计数器的维护、读数据多路选择，并在配置表写入时与收发核心执行许可握手。模块输出全部来自寄存器，结构见图~\ref{fig:cfg}。

\section{地址映射}

配置地址为 18 位字地址，映射见表~\ref{tab:addrmap}。各条 VL 通过自己的 \code{en} 位控制启用。对表外地址的写请求予以忽略，读请求返回 0。

\begin{longtable}{>{\raggedright\arraybackslash}p{4.2cm}>{\centering\arraybackslash}p{1.2cm}>{\raggedright\arraybackslash}p{\dimexpr\linewidth-5.4cm-6\tabcolsep\relax}}
\caption{配置地址映射}\label{tab:addrmap}\\
\longtablehead{3}{地址 & 读写 & 内容}
\code{0x00000} & 只读 & \code{VERSION}，\code{0xC3110000} \\
\code{0x00001} & 只读 & \code{STATUS}，位 0 为 \code{tx_ready} \\
\code{0x00002} & 读写 & \code{USER_ID[15:0]}，填入源 MAC 第 9--10 字节 \\
\code{0x00040}--\code{0x00046} & 只读 & 发送计数器：\code{accept}、\code{drop_dis}、\code{drop_full}、\code{drop_long}、\code{drop_nopage}、\code{tx_a}、\code{tx_b} \\
\code{0x00048}--\code{0x0004F} & 只读 & A 网接收计数器，对应 \code{ev_a} 位 0--7 \\
\code{0x00050}--\code{0x00057} & 只读 & B 网接收计数器，对应 \code{ev_b} 位 0--7 \\
\code{0x00058}、\code{0x00059} & 只读 & \code{rm_drop}、\code{deliver} \\
\code{0x0005A}、\code{0x0005B} & 只读 & \code{a_gdrop}、\code{b_gdrop}（保留，恒为 0） \\
\code{0x00080}--\code{0x000BB} & --- & RTC、SNMP、615A 寄存器（附录表~\ref{tab:regs}） \\
\code{0x01000} $+$ VL$\times$8 $+$ 字号 & 只写 & \code{tx_cfg} \\
\code{0x10000} $+$ VLID & 只写 & \code{rx_map}，数据取位 12:0 \\
\code{0x20000} $+$ idx & 只写 & \code{rx_cfg}，idx $<$ 4096 \\
\end{longtable}

三张表按地址高位译码：地址位 17:10 为 \code{0000_0100} 时选中 \code{tx_cfg}，位 17:16 为 \code{01} 时选中 \code{rx_map}，位 17:12 为 \code{10_0000} 时选中 \code{rx_cfg}。

\section{配置表格式}

三张配置表的字段见表~\ref{tab:cfgfmt}。复位后三张表全部清零，即全部 VL 未启用、全部 VLID 未映射。

\begin{longtable}{>{\raggedright\arraybackslash}p{2.6cm}>{\centering\arraybackslash}p{1.4cm}>{\raggedright\arraybackslash}p{2.6cm}>{\raggedright\arraybackslash}p{\dimexpr\linewidth-6.6cm-8\tabcolsep\relax}}
\caption{配置表字段}\label{tab:cfgfmt}\\
\longtablehead{4}{表与字 & 位 & 字段 & 含义}
\code{tx_cfg} 字 0 & 0 & \code{en} & VL 启用 \\
 & 1、2 & \code{net_a}、\code{net_b} & 从 A、B 网发送 \\
 & 8:3 & \code{qmax} & 队列最多排队的消息数，0--63 \\
 & 19:9 & \code{max_payload} & 单帧最大有效载荷 $L_\text{max}$，1--1471 字节（应用消息 $M>L_\text{max}$ 且 $L_\text{max}\ge 8$ 时触发 IPv4 分片，消息总长上限 $\MaxMessageBytes$~B） \\
\code{tx_cfg} 字 1 & 31:0 & \code{bag} & BAG，单位为拍，$1~\text{ms}=125000$，不超过 $2^{30}$ \\
\code{tx_cfg} 字 2 & 15:0 & VLID & 目的 MAC 低 16 位（顶层 \code{AFDX_top} 自动窥探同步至 \code{afdx_app_tx}） \\
\code{tx_cfg} 字 3、4 & 31:0 & 源 IP、目的 IP & IPv4 地址 \\
\code{tx_cfg} 字 5 & 31:16、15:0 & 源端口、目的端口 & 消息不带端口时使用的默认 UDP 端口 \\
\code{tx_cfg} 字 6、7 & --- & 保留 & 写 0 \\
\midrule
\code{rx_map}[VLID] & 12 & \code{valid} & 该 VLID 已映射（顶层 \code{AFDX_top} 自动窥探同步至 \code{afdx_app_rx}） \\
 & 11:0 & \code{idx} & 接收 VL 索引，0--4095 \\
\midrule
\code{rx_cfg}[idx] & 0 & \code{en} & VL 启用 \\
 & 1、2 & \code{net_a}、\code{net_b} & 接收 A、B 网的帧 \\
 & 3、4 & \code{ic_en}、\code{rm_en} & 启用完整性检查、冗余管理 \\
 & 16:5 & \code{skew_max} & 两网最大到达时间差，单位 32.768~µs \\
 & 27:17 & \code{max_payload} & 单帧最大载荷 $L_\text{max}$，1--1471 字节（后续 IPv4 分片允许 $\le L_\text{max}+8$） \\
\end{longtable}

\section{配置读写与仲裁}

读数据的多路选择分两拍完成，读延迟固定为 2 拍：第 1 拍按地址低 3 位从四组数据源（全局寄存器、发送计数器、A 网计数器、B 网计数器）及三项服务寄存器中并行预选出候选字并锁存，同时按地址位 17:3 锁存组号；第 2 拍按组号从各路寄存值中选出目标字送至 \code{cfg_rdata}，并拉高 \code{cfg_rvalid}，每拍至多经过一级 8 选 1 逻辑。

\figN[配置模块]{f30_cfg}{配置事务示例：(a) 暂存对 \code{rx_cfg[7]} 的写入，等待许可后写 RAM，并对比 RX/TX 的就绪重武装；(b) 读回 \code{USER_ID}=42 的两拍候选锁存与选择}{fig:cfg}

\code{cfg_ready} 为 1 时，总线在当拍接受 \code{cfg_wr} 写请求：普通寄存器和服务寄存器直接写入；配置表写请求则先装入暂存寄存器并拉低 \code{cfg_ready}，等待对应核心给出 \code{cfg_ok}（第~\ref{sec:intif}~节）。\code{tx_ready} 为 0 时，配置表写暂存寄存器保持清零，\code{cfg_ready} 为 0，不接受表初始化期间的写请求。

SNMP 代理读写管理对象使用同一条总线。顶层的仲裁规则是：主机在本拍写入，或主机与代理读取不同地址时，优先处理主机请求；主机读与代理读同一地址时合并为一次读，两者各自在 2 拍后取得数据；代理写只在主机无读写且 \code{cfg_ready} 为 1 时授予总线访问权。主机与代理的读数据分别由各自的 2 拍有效标记对齐。

\section{配置写与查表的冲突仲裁}\label{sec:cfgok}

配置表的读口归收发引擎，写口归配置总线。配置写通过 \code{cfg_ok} 与引擎读请求互斥，避免同步 RAM 同拍同址读写冲突。各引擎的许可条件为：
\begin{align}
\text{发送侧：}\ & \text{cfg\_ok} = \text{enable} \wedge (\text{op}=\text{IDLE}) \wedge \neg\text{want\_enq} \wedge \neg\text{want\_grant} \\
\text{接收侧：}\ & \begin{aligned}[t]
  \text{cfg\_ok} = \text{enable} &\wedge (\text{op}\in\{\text{IDLE},\text{AGE}\}) \\
  &\wedge \neg\text{desc\_a\_valid} \wedge \neg\text{desc\_b\_valid}
\end{aligned}
\end{align}
\code{cfg_ok} 由引擎自身的状态寄存器与同级模块的寄存输出（入队请求、帧描述有效、\code{grant_ready}）组合生成。\code{afdx_cfg} 接收一次写表请求后，保持写使能（\code{tx_cfg_we}、\code{rx_map_we} 或 \code{rx_cfg_we}）及地址、数据；当对应核心的 \code{cfg_ok} 为 1 时，引擎本拍不置位对应配置表的读使能寄存器，核心模块在该拍把写请求锁存到写口寄存器并于下一拍写入 RAM，\code{afdx_cfg} 撤销写请求。接收侧在许可后的下一拍恢复 \code{cfg_ready}；发送侧通过 \code{tx_cfg_ok_q} 与 \code{tx_cfg_rearm} 再延后一拍恢复就绪，保证本地应答与真实写请求完成对齐。时序为：
\begin{equation}
\begin{aligned}
&\text{cfg\_ready}(t+1)=0,\qquad
u=\min\{\,\tau \ge t+1 : \text{cfg\_ok}(\tau)=1\,\},\\
&\text{cfg\_ready}_{\mathrm{RX}}(u+1)=1,\qquad
\text{cfg\_ready}_{\mathrm{TX}}(u+2)=1.
\end{aligned}
\end{equation}
其中 $t$ 为总线接受写请求的时刻，$u-t$ 为写表等待时间，上界为 64 拍（第~\ref{sec:budget}~节）。


\section{复位与初始化}

复位与表初始化过程见第~\ref{sec:clkrst}~节。配置模块复位后，各寄存器恢复默认值，27 个计数器清零。\code{tx_ready} 为 0 期间，配置表写暂存寄存器保持清零，\code{cfg_ready} 保持为 0；此时不接受配置写请求，读访问照常进行。

\section{统计计数器}

27 个计数器均为 32 位，计满后回绕至 0，仅在复位时清零。为隔离发送核心与接收核心长距离跨模块引入的时序关键路径，\code{afdx_cfg} 对来自 TX/RX 通路的 5 路事件统计脉冲（\code{ev_tx}、\code{ev_a}、\code{ev_b}、\code{ev_rm_drop}、\code{ev_deliver}）插入了\textbf{两级流水寄存}（\code{_p} 与 \code{_q}）后再驱动 32 位加法器，降低了跨芯片远距离脉冲直接驱动大位宽累加器的布线延迟。

计入事件输入两级流水打拍（\code{_p} 与 \code{_q}）及读数据通路的两级流水结构，在 $t$ 拍发出配置总线读请求，$t+2$ 拍返回对应事件脉冲截至 $t-3$ 拍在核心源头产生的累计次数（对应配置模块内部计数值在 $t-1$ 拍及以前的累加状态）。计数器分为：
\begin{itemize}
  \item \textbf{发送（7 个）：}接纳的消息数；因 VL 未启用、队列满、超长、无空闲页而丢弃的消息数；A、B 网发出帧数。
  \item \textbf{接收，每网 8 个：}正确帧、FCS 错误、帧长错误、帧头错误、未知 VL、配置不允许、完整性检查丢弃、溢出。
  \item \textbf{接收，两网合计 2 个：}冗余管理丢弃、交付。
  \item \textbf{保留 2 个：}\code{a_gdrop}、\code{b_gdrop}，地址保留，恒为 0。
\end{itemize}

\section{异常处理汇总}\label{sec:errsum}

表~\ref{tab:errsum} 汇总全设计对各类异常的检测位置、处理方式和统计。发送侧拒收消息后，先回收已分配的页，再继续接收并丢弃剩余字节。

\begin{longtable}{>{\raggedright\arraybackslash}p{3.6cm}>{\raggedright\arraybackslash}p{2.0cm}>{\raggedright\arraybackslash}p{5.0cm}>{\raggedright\arraybackslash}p{2.8cm}}
\caption{异常处理汇总}\label{tab:errsum}\\
\longtablehead{4}{异常 & 检测位置 & 处理 & 统计}
\multicolumn{4}{l}{\textit{发送与顶层发送包装}} \\*
顶层发送 VLID 未映射 & \code{afdx_app_tx} & 整条消息丢弃，不进入核心 & \code{ev_app[0]} \\
顶层发送超 $\MaxMessageBytes$~B 或包装缓存满 & \code{afdx_app_tx} & 回滚包装写指针，丢弃本条消息 & \code{ev_app[1]} \\
VL 或两网均未启用 & \code{tx_engine} & 拒绝入队，回滚回收已用缓存页 & \code{drop_dis} \\
队列已满 & \code{tx_engine} & 同上 & \code{drop_full} \\
核心消息超过 $\MaxMessageBytes$~B，或分片时 $L_\text{max}<8$，或单帧超过 $L_\text{max}$（当配置不分片） & \code{tx_engine}、\code{tx_buf} & 从回滚 RAM 回收全部已分配页，继续接收并丢弃剩余字节 & \code{drop_long} \\
无空闲页 & \code{tx_buf} & 同上 & \code{drop_nopage} \\
MAC 发送下溢（TX Underrun） & \code{mac_tx} & 发送未取反的 FCS，使对端判为错误帧，继续接收并丢弃本帧剩余字节 & 不计 \code{tx_a}/\code{tx_b} \\
\midrule
\multicolumn{4}{l}{\textit{接收、分片重组与顶层接收包装}} \\*
前导码错误、帧长越界 & \code{mac_rx} & 整帧丢弃 & 帧长错误 \\
FCS 错误 & \code{mac_rx} & 整帧丢弃 & FCS 错误 \\
固定字段、长度、非 8B 对齐分片或校验和错误 & \code{rx_parse} & 回退预写入指针 & 帧头错误 \\
环形缓冲剩余空间不足 & \code{rx_parse} & 不写入本帧 & 溢出（\code{ev_a/b[7]}） \\
VLID 未映射 & \code{rx_engine} & 回退预写入指针 & 未知 VL \\
VL 或本网未启用、单帧超 $L_\text{max}$（后续片超 $L_\text{max}+8$） & \code{rx_engine} & 回退预写入指针 & 配置不允许 \\
SN 不连续 & \code{rx_engine} & 回退预写入指针，PSN 仍更新为本次 SN & 完整性检查丢弃 \\
另一网已交付的副本 & \code{rx_engine} & 回退预写入指针 & \code{rm_drop} \\
交付队列已满，或分片缺首片/乱序/重叠/16 槽全满 & \code{rx_engine}、\code{rx_reasm} & 回退预写入指针，拒绝进入重组槽位 & 溢出/拒绝（\code{ev_a/b[7]}） \\
末片总长不符、非零全包 UDP 校验错或 2~s 老化超时 & \code{rx_reassembly} & 废弃重组槽位并释放上下文 & 溢出/拒绝（\code{ev_a/b[7]}） \\
顶层接收包装缓存溢出 & \code{afdx_app_rx} & 整包丢弃，保护已入队消息边界 & \code{ev_app[2]} \\
表初始化完成前开始的帧 & \code{rx_parse} & 整帧忽略 & 不计 \\
\midrule
\multicolumn{4}{l}{\textit{服务与配置}} \\*
SNMP community 不符 & \code{afdx_snmp} & 不响应；Trap 启用时发 \code{authenticationFailure} & \code{BAD_COMMUNITY} \\
SNMP 请求超过 256 字节或代理忙 & \code{afdx_snmp} & 端口匹配后丢弃，由管理端超时重发 & 忙时计入 \code{REQUESTS} \\
615A 数据/文件检查错、状态传输失败、请求错误 & \code{afdx_615a} & 数据/文件错：ERROR 及失败状态文件；状态传输失败（含重试耗尽）：直接结束；名称/模式错：仅 ERROR；LCI/LNO：初始化拒绝 & \code{STATUS} 原因码 1--5（第~\ref{sec:a615}~节） \\
SNTP 响应版本、模式或 Originate 不符 & \code{afdx_rtc} & 丢弃，下一周期重发请求 & --- \\
表初始化期间写配置 & \code{afdx_cfg} & \code{cfg_ready}$=0$，不接受 & --- \\
未映射的地址 & \code{afdx_cfg} & 忽略写请求，读出 0 & --- \\
\end{longtable}

\section{主机编程流程}\label{sec:prog}

\subsection{初始化与启用}

主机经配置总线按下列顺序使端系统进入工作状态。每次写表都要等 \code{cfg_ready} 为 1（第~\ref{ch:if}~章）。
\begin{enumerate}
  \item \textbf{初始化：}释放复位后轮询 \code{STATUS}，等待位 0（\code{tx_ready}）变为 1，约 524~µs。
  \item \textbf{端系统标识：}写 \code{USER_ID}。
  \item \textbf{发送 VL：}对每条发送 VL 先写字 1--5（BAG、VLID、IP、默认端口），最后写字 0，同时设置 \code{en}、网络选择、\code{qmax} 和 \code{max_payload}。
  \item \textbf{接收 VL：}对每条接收 VL 先写 \code{rx_cfg}[idx]，再写 \code{rx_map}[VLID]$=\{1, \text{idx}\}$，使映射生效时配置已经就绪。
  \item \textbf{服务：}为每项服务配置一条发送 VL（目的 IP、端口与服务端口一致），写服务的 \code{TX_VL}、端口和 community 等寄存器，最后写 \code{ENABLE}。
  \item \textbf{开始收发：}应用开始提交消息。VL 较多时按第~\ref{sec:jitter}~节错开各 VL 首条消息的提交时刻。
\end{enumerate}

\subsection{运行中更新与服务配置}

\code{tx_cfg} 的 \code{en} 仅在入队时检查；已排队的消息继续参与授权，授权阶段逐拍读取字 0--5。修改一条发送 VL 时，主机应先停止向它提交消息，再把字 0 的 \code{en} 写 0，并保持该字其余字段不变；待该 VL 已提交的消息全部发出后，修改其余字，最后重新启用。

通过 \code{cfg_ok} 保证配置写与表读互斥，多字配置按字依次生效。接收侧的 \code{rx_cfg} 和 \code{rx_map} 各自通过单拍整字写入更新；两表的写入分别生效，历史状态在配置更新后保留。

应用与服务的配置还须满足以下条件：
\begin{itemize}
  \item 应用须等待表示收发核心均已就绪的 \code{tx_ready} 有效后，再提交消息。\code{app_tx_ready} 由应用仲裁选中条件与发送核心的 \code{m_ready} 相与生成。
  \item 各服务的 UDP 端口须互不冲突；RTC 的模式与时间，以及 SNMP 的 \code{ENABLE}、\mbox{community} 和 \code{TX_VL}，仅在对应服务没有在途消息时修改。
  \item 615A 忙时，忽略对启用、发送 VL、请求端口、超时、重传上限和文件版本的控制写请求；写 \code{0xB4} 后至少等待 4 拍，且只在非忙状态下读 \code{0xB5}。
\end{itemize}

\section{配置约束与网络规划规格}\label{sec:cfgrule}

主机须检查各项配置的组合是否满足带宽、抖动、缓存和分片重组并发度要求。结合官方 1030 答疑 R138-3/R139-3，本设计声明的\textbf{端系统满载网络规划规格}为单端 63 条发送 VL / 63 条接收 VL（$\text{BAG}=1$~ms，$L_\text{max}=1471$~B，双网冗余全双工，物理线路占用率 $77.5\%$，留有约 $22.5\%$ 突发排空与管理服务裕量）。具体约束及超限处理如下：
\begin{itemize}
  \item \textbf{带宽：}A、B 两网由同一个 MAC 状态机同拍驱动，只发 A 网或只发 B 网的帧同样占用一个发送时隙，所以约束针对全部发送 VL，不分网络：$\sum_v C(n_v)/\text{BAG}_v \le 1$，其中 $n_v$ 为单帧最大载荷 $L_\text{max}$，$C(n_v)$ 为第~\ref{sec:budget}~节式~(3-1) 的线路占用拍数，服务 VL 计入求和。长包分片消息在每个 BAG 周期仅调度 1 个分片帧，单周期带宽占用与单帧 VL 一致。
  \item \textbf{抖动：}BAG 取 $2^k$~ms；周期业务宜错开各 VL 首条消息相位，在 63~VL 满载规格下即使发生旋转三元组背靠背突发（B2），同 VL 最大端到端极差抖动仍保持在 \JitSixtyThreeBTwo~µs（$<100$~µs），而超规格 81~VL 对齐突发将因排队逼近 1~ms 周期而超出百微秒门限。
  \item \textbf{发送与重组缓存：}发送核心共 256 页（128~KiB），单帧 $1471$~B 占 3 页，$8192$~B 最大长包占 16 页；主机规划在途驻留页数须满足 $\sum_v \text{qmax}_v \times \lceil M_v / 512 \rceil + 9 \le 256$（例如 60 条短包加 $1472/3000/8192$~B 三条长包共占 205 页，安全满足；多条并发 $8192$~B 长包则须受 256 页总预算与顶层 32~KiB 包装缓冲约束）。接收侧配备 \ReasmContexts{} 个并发重组槽位（奇偶 VL 索引各 8 个槽位），同一时刻处于跨 BAG 分片重组中的活跃长包 VL 数在奇偶各 Bank 内均不得超过 8 条，且同 VL 新长包须在前一长包末分片发出后再提交。
  \item \textbf{交付：}交付口每拍一字节。若某条 VL 关闭冗余管理，且同时接收两网数据，则每条消息交付两次。所有接收 VL 的载荷交付总速率须低于 1 字节每拍；两网独立满速且都要交付时，所需交付速率超过该上限，超出部分按溢出计数。
  \item \textbf{顺序与冗余窗口：}跨 VL 的发送顺序由各 VL 的 BAG 与候选选择共同决定；接收分片必须在同一 VL 内按递增片偏移顺序到达，乱序分片将被拒绝。接收副本的到达间隔超过 \code{skew_max} 时，冗余记录按过期处理。
  \item \textbf{取值范围：}\code{bag} $\le 2^{30}$ 拍；\code{qmax} $\le 63$；\code{max_payload} 为 1--1471（启用分片时 $\ge 8$）；应用消息总长 $1\le M\le \MaxMessageBytes$~B；\code{skew_max} 小于过期清理门限 \code{0x4000} 个计时单位（12 位字段自然满足）。
\end{itemize}

